\item[\textbf{Problem 5}]
Assume that the AN's classification of each of the points above is correct.
\textbf{List} a new labeled data item which, if added to the data set above,
would necessarily cause the AN to misclassify at least one data item. (That is,
given this new data item, there would be no set of possible weights that would
allow this AN to correctly classify all data items.) \textbf{Explain} why the AN
would necessarily misclassify at least one data item.

\Solution

After inspecting the plotted data points, we can see that the existing points
form a cluster of \textbf{Class 1} items. To solve the problem, we need to
introduce a new point that breaks this pattern, making it impossible to draw a
single straight line that separates two classes. A reliable method is to choose
a point that lies on the line segment between two existing points. I selected
the existing points $(0.0, 0.0)$ and $(1.0, 0.2)$. If we place a new
\textbf{Class 0} point at their midpoint, $(0.5, 0.1)$, we create a conflict.
This arrangement results in two Class 1 points flanking a Class 0 point on the
same line. No single decision boundary can isolate the middle point from the
outer two without misclassifying at least one of them.

\begin{minipage}{\linewidth}
    \centering
    \includegraphics[width=0.8\linewidth]{problem-4-graph-with-class-0.pdf}
    \captionof{figure}{Decision boundary with the Class 1 cluster and an outlier within it}
    \label{figure:slope-graph-with-classified-points-and-an-outlier}
\end{minipage}
